\documentclass[a4paper,10pt]{article}
\usepackage[utf8]{inputenc}
\usepackage[brazil]{babel}
\usepackage{graphicx}
\usepackage{makeidx}
\usepackage{spverbatim}
\makeindex


%opening
\date{\today}
\begin{document}
    \begin{titlepage}
        \begin{figure}
            \centering
          \includegraphics[width=7cm,keepaspectratio=true]{imagens/unicamp.png}

        \end{figure}
        \begin{center}
            \huge{Universidade Estadual de Campinas}

        \vfill
        \textbf{\LARGE{Manual}}
        \vfill
        \end{center}

        \begin{flushleft}
            \begin{tabbing}
                Autores\qquad\qquad\= Gabriel Bueno de Oliveira 139455 \\
                    \>Joao Guilherme Daros Fidelis 136242 \\
                    \>Lucas Henrique Morais 136640 \\
                    \>Matheus Yokoyama Figueiredo 137036\\
                    \>Pedro Rodrigues Grijo 139715\\
            \end{tabbing}
        \end{flushleft}
    \end{titlepage}

\newpage
\section{Setup do Cluster}
    \subsection{Acesso}
    O acesso ao cluster pode ser feito de fora e de dentro do IC (Instituto de Computação - UNICAMP).
        \subsubsection{Interno ao IC}
        Para acceso interno ao ic, basta acessar por ssh o endereço cbn6.lab.ic.unicamp.br. No linux, isso pode ser feito através de:
        \begin{spverbatim}
        # ssh cbn6.lab.ic.unicamp.br
        \end{spverbatim}

        Após a conexão ser feita, insira sua senha.
        \subsubsection{Externo ao IC}
        Para o acesso externo, é necessário fazer uma conexão de duas fases ou usar um tunnel por uma máquina do ic. Para acesso remoto a uma máquina do ic, no linux, utilize \textbf{um} dentre os seguintes comandos:
        \begin{spverbatim}
        # ssh -l raXXXXXX ssh.students.ic.unicamp.br
        # ssh -l raXXXXXX ssh2.students.ic.unicamp.br
        \end{spverbatim}

        Após a conexão ser feita, insira sua senha. Após feita a credencial, prossiga do mesmo modo como se tivesse acessando de dentro do IC.

\section{Instalação e Configuração}
    \subsection{PostgreSQL}
        \subsubsection{Instalação}
        A versão utilizada do PostgreSQL é a versão 9.5, que pode ser obtida seguindo os seguintes comandos:
        \begin{spverbatim}
        # sudo su
        # cat /etc/apt/sources.list.d/postgresql.list
        # deb http://apt.postgresql.org/pub/repos/apt/ trusty-pgdg main
        # sudo apt-get install wget ca-certificates
        # wget --quiet -O - https://www.postgresql.org/media/keys/ACCC4CF8.asc | sudo apt-key add -
        # sudo apt-get update
        # sudo apt-get upgrade
        # sudo apt-get install postgresql-9.5

        \end{spverbatim}

        \subsubsection{Configuração}
        Para configurar o acesso ao banco de dados deve-se alterar o arquivo de configuração principal. Nele são definidos toda a configuração de acesso ao banco de dados e os diretórios de dados que são respectivamente:

        \begin{spverbatim}
        hba_file = '/etc/postgresql/9.5/main/pg_hba.conf'
        data_directory = '/var/lib/postgresql/9.5/main'
        \end{spverbatim}

        O arquivo pg\_hba.conf define regras para a conexão ao banco. Por padrão, são só permitidas conexões locais: o usuário administrador do banco (postgres, quando instalado por gerenciador de pacotes) tem acesso via socket e sem senha. os demais usuários, via interface local e com senha:
        \begin{spverbatim}
        local   all         postgres            peer

        host    all         all         127.0.0.1/32            md5
        \end{spverbatim}

    Para criar o primeiro usuário siga os seguintes comandos:
        \begin{spverbatim}
            # sudo su
            # su postgres
            # psql
            =# CREATE ROLE <usr_name> with superuser createdb createrole login;
            =# ALTER ROLE <usr_name> PASSWORD ‘<password>’;
        \end{spverbatim}

        \subsubsection{Carregando o Dump}
        Para carregar o ddump no banco deve-se criar o banco de dados (tpcw) e um usuário tpcw-user.
        \begin{spverbatim}
        # pg_restore -d tpcw tpcw-database-dump.tar
        \end{spverbatim}


    \subsection{Tomcat}
        \subsubsection{Instalação}
        Para instalar o Tomcat (usaremos o tomcat7), precisamos instalar o pacote do tomcat e suas dependencias. Para isso, no linux, podemos usar o apt-get:
        \begin{spverbatim}
        # sudo apt-get install tomcat7
        \end{spverbatim}

        A instalação irá criar um usuário tomcat7 dentro de um grupo, tomcat7. Os arquivos estão em \$CATALINA\_BASE, que na instalação do Ubuntu é /var/lib/tomcat7.

        Para teste, acesse em uma máquina do IC o cbn6.lab.ic.unicamp.br:8080.

        \subsubsection{Integração com o PostgreSQL}
        Para integração com o PostgreSQL é necessário fazer o download do .jar do PostgreSQL JDBC para ser adicionado a pasta \$CATALINA\_BASE/lib. Após fazer o download, modifique os seguintes arquivos com o conteúdo indicado:

        \begin{itemize}
        \setlength{\itemindent}{-.0in}
            \item \$CATALINA\_BASE/conf/context.xml: http://www.ggte.unicamp.br/eam/pluginfile.php/188108/
            mod\_wiki/attachments/132/tomcat7-context.xml
            \item \$CATALINA\_BASE/conf/server.xml: http://www.ggte.unicamp.br/eam/pluginfile.php/188108/
            mod\_wiki/attachments/132/tomcat7-server.xml. Para este, modifique os campos de username, password e possivelmente url.
        \end{itemize}

        Após essas configurações, reinicie o tomcat:

        \begin{spverbatim}
        # invoke-rc.d tomcat7 restart
        \end{spverbatim}

        \subsection{TPC-W}

        As aplicações implantadas no Tomcat são configuradas em \$CATALINA\_BASE/webapps. Para ter permissão de criar aplicações nesta pasta, o usuário tem que ser do grupo tomcat7, para fazer isso, no linux:

        \begin{spverbatim}
        # adduser [usuário_comum] tomcat7
        \end{spverbatim}

        Para instalação default do TPC-W no tomcat7, crie a seguinte árvore:
        \begin{spverbatim}

            \$CATALINA_BASE/webapps

                tpcw

                     Images

                     WEB-INF

                        classes

                        web.xml


        \end{spverbatim}

        Na pasta tpcw, rode os seguintes comandos:
        \begin{spverbatim}
            # sudo cp /nfs/servlet/Images.tar.gz .
            # sudo tar -xzf Images.tar.gz
            # sudo cp /nfs/servlet/tpcw-servlets.tar.gz
            # sudo tar -xzf tpcw-servlets.tar.gz
            # cd tpcw-servlets.tar.gz
            # sudo make
        \end{spverbatim}

        No fim, copie todas as *.class geradas para a pasta WEB-INF/classes e copie o web.xml para pasta WEB-INF/. Reinicie o Tomcat.

       \subsection{RBE}
        O RBE (Remote Browser Emulator) é o programa empregado para gerar carga para o TPC-W. É um aplicativo que emula um conjunto de clientes que acessam o lado servidor do TPC-W, que implementa uma loja de livros. Os clientes operam como se fossem navegadores web.
        O código fonte do RBE, assim como arquivo README e o script em Python que realiza a análise das saídas estão disponíves na pasta do RBE.

        \subsubsection{Automatização}
        Para facilitar a coleta dos dados foi implementado um script em bash que executa o rbe com variando os parâmetros e as saídas, já executa o script analyse.py e gera os gráficos.
        Para executar o script basta executar:
        \begin{spverbatim}
        # ./script.sh
        \end{spverbatim}

\section{Entrando e conhecendo as novas máquinas}
    Obtivemos acesso a duas máquinas diferentes do cluster que contém uma instância do postgres cada, assim, uma máquina conterá o banco primário e a outra o secundário.\\
    O banco primário fica na máquina dbmaster2. Essa máquina tem IP 10.1.2.10 e o banco fica na porta 5434.\\
    Já o banco secundário fica na máquina dbslave2. Essa máquina tem IP 10.1.2.20 e o banco fica na porta 5434.

    Para acessar o dbmaster2, basta enquanto logado no cluster cbn6, dar ssh no seu endereço, assim:
    \begin{spverbatim}
    # ssh dbmaster2
    \end{spverbatim}

    Para acessar o dbslave2, basta enquanto logado no cluster cbn6, dar ssh no seu endereço, assim:
    \begin{spverbatim}
    # ssh dbslave2
    \end{spverbatim}

    As máquinas já vieram com a instalação do Postgres 9.5.1.

    As pastas de dados para os bancos de dados em ambas as máquinas fica em \verb|/var/lib/postgresql/9.5/grupo06/|.

    As pastas de configurações (como postgres.conf e pg\_hba.conf), ficam em \verb|/etc/postgresql/9.5/grupo06/|.

    Também é necessário configurar e iniciar um banco de dados inicial no banco primário em dbmaster2.

    Faça login na máquina dbmaster2 e crie um usuário "tpcw-user" como é ensinado na seção 2.1.2 e faça uma restauração do dump como na seção 2.1.3.

\section{Instalando a Replicação}
    Era necessário instalar o serviço de Streaming Replication do banco primário para o secundário, para que mudanças feitas no banco primário sejam copiadas para o secundário. Para isto, seguimos tutoriais sugeridos pelo professor.
    Agora, temos dois bancos de dados em duas máquinas diferentes.
    \subsection{Editar configurações no mestre}
      Na máquina do banco primário, dbmaster2, vamos criar um usuário "replication" no postgres, que é um usuário com privilégios de replicação. O usuário terá a senha 123456. Para isso, execute os comandos:
      \begin{spverbatim}
      # sudo su
      # su postgres
      # psql
      # CREATE ROLE replication WITH REPLICATION PASSWORD ‘123456’ LOGIN
      \end{spverbatim}

      Agora, vamos editar o arquivo /etc/postgresql/9.5/grupo06/postgresql.conf da máquina dbmaster2.
      Na seção \textbf{CONNECTIONS AND AUTHENTICATION}, adicione a linha:
      \begin{spverbatim}
      listen_addresses = '*'
      \end{spverbatim}

      Na seção \textbf{WRITE AHEAD LOG}, adicione as linhas:
      \begin{spverbatim}
      wal_level = hot_standby
      \end{spverbatim}

      Em \textbf{REPLICATION}, coloque as linhas:
      \begin{spverbatim}
      max_wal_senders = 5
      wal_keep_segments = 32
      \end{spverbatim}

      Finalmente, em \textbf{ARCHIVING}, adicionamos as linhas:
      \begin{spverbatim}
      archive_mode = on
      archive_command = 'test ! -f /var/lib/postgresql/9.5/grupo06/archivedir/\%f && cp \%p /var/lib/postgresql/9.5/grupo06/archivedir/\%f'
      \end{spverbatim}

     Depois, crie a pasta do comando acima com:
      \begin{spverbatim}
      # mkdir /var/lib/postgresql/9.5/grupo06/archivedir/
      \end{spverbatim}

      Veja que a pasta /var/lib/postgresql/9.5/grupo06/archivedir/ será criada e nela ficará contido todos os arquivos log gerados pelos bancos de dados.
      Agora, temos que editar o arquivo /etc/postgresql/9.5/grupo06/pg\_hba.conf. Adicione a linha no final:
      \begin{spverbatim}
      host    replication     replication     10.1.2.20/32    md5
      \end{spverbatim}

      Essa linha permite que o banco em standby consiga se conectar com o primário para fazer a replicação.
    \subsection{Fazer base backup}
        Devemos fazer um base backup do servidor mestre para o escravo. Utilizaremos o comando pg\_basebackup.\\
        Para isso, faça login na máquina do escravo, a dbslave2.\\
        Primeiro, a pasta de dados do banco de dados que receberá o base backup deve estar vazia. Logo, apagaremos tudo dentro dela. Depois, utilizaremos o comando de base backup para se conectar no banco de dados mestre e fazer a cópia.\\
        Rode o comando:
        \begin{spverbatim}
        rm /var/lib/postgresql/9.5/grupo06/* -r
        pg_basebackup -h 10.1.2.10 -p 5434 -D /var/lib/postgresql/9.5/grupo06 -P -U replication --xlog-method=stream -R
        \end{spverbatim}

        Ao ser pedido uma senha, digite 123456.\\
        A flag -R no final, faz com que o arquivo recovery.conf seja gerado automaticamente no processo.

    \subsection{Configurar o slave}
        Na máquina escrava, edite o arquivo \verb|/etc/postgresql/9.5/grupo06/postgresql.conf|, adicionando as mesmas linhas que fez na máquina mestre, como nos passos anteriores. Isso é necessário pois quando o primário falhar, o secundário deve se comportar exatamente como o primário. Entretanto, é necéssario adicionar mais uma linha, na seção \textbf{STANDBY SERVERS}, adicione a linha:
        \begin{spverbatim}
        hot_standby = on
        \end{spverbatim}

        Isso faz com que o banco escravo seja read-only e aceite apenas comandos de leitura.

\section{Configuração do HAProxy}

    \subsection{Instalação}

        Para instalar o HAProxy, tivemos apenas que adicionar ao nosso ambiente de execução o repositório apt-get relacionado e rodar o comando apt-get install apropriado, como segue:

        \begin{spverbatim}
        # sudo add-apt-repository ppa:vbernat/haproxy-1.6
        # sudo apt-get update
        # sudo apt-get install haproxy
        \end{spverbatim}

        Podemos verificar que a versão correta foi então instalada com o seguinte comando:

       \begin{spverbatim}
        # haproxy -v
        \end{spverbatim}

    \subsection{Configuração}

        Uma vez instalado o HAPROXY, temos ainda de configurá-lo de modo a:

        \begin{enumerate}
        \item Informá-lo da existência de dois endereços de destino distintos
        \item Permitir que seu backend escolha o endereço de destino adequado de acordo com o resultado de um health-check periódico de cada um dos alvos
        \item //Requisitar que o HAProxy escreva mensagens de log relevantes em /var/log
        \end{enumerate}

        Para tal, é necessário criar em \verb|/etc/haproxy| um arquivo haproxy.cfg com parâmetros adequados.
        A seguir, comentamos o conteúdo de um arquivo mínimo desse tipo que é capaz de prover o requerido em (1) e (2):

       \begin{spverbatim}
        listen psql

       bind    *:4002

       mode  tcp

       option  pgsql-check user haproxy

       server  replicaOne  10.1.2.10:5434 check

       server  replicaTwo  10.1.2.20:5434 check backup
        \end{spverbatim}

        A linha 02 garante que toda a comunicação ocorrendo através da porta 4002 deverá ser interceptada.

        A linha 03 indica  que o haproxy deverá funcionar como um tcp proxy – isto é, transparentemente interceptando e redirecionando pacotes do tipo TCP.

        A linha 04 configura o haproxy para fazer contatos de health-check com os bancos replicados usando o usuário haproxy.

        As linhas 05 e 06 definem as réplicas do banco como os dois destinos possíveis para os pacotes TCP interceptados. Além disso, a opção “backup” da linha 06 garante que o servidor replicaTwo não seja contatado antes que replicaOne falhe.

        Por fim, as opções “check” ocorrendo nas duas linhas indicam que haproxy deverá verificar periodicamente a disponibilidade dos dois bancos, alternando o destino dos pacotes interceptados de acordo com essa avaliação.

    \subsection{Startup, restart e shutdown}

        Para realizar o startup, restart e shutdown do HAProxy, respectivamente, basta usar o comando “service” do Linux, como segue:

        \begin{spverbatim}
        sudo service haproxy start
        sudo service haproxy restart
        sudo service haproxy stop
        \end{spverbatim}


\section{Automatização da coleta de dados}

            Todos os scripts mencionados nessa seção estão na pasta \verb|scripts|.

	A segunda parte do experimento foi dividida em 3 partes. Para executar os testes da primeira parte, foram gerados os scripts \verb|arc2_exp1_script.sh| e \verb|arc2_exp_analise.sh| que devem ser executados da seguinte forma:

 	\begin{spverbatim}
	# ./arc2_exp1_script.sh
	# ./arc2_exp_analise.sh 1
	\end{spverbatim}
\noindent
\newline
O argumento passado para o segundo script é simplesmente para geração de arquivos com o nome desejado.
\newline

Para executar os testes da segunda parte, foi gerado o script \verb|arc2_exp2_script.sh|. Seguem os comandos para executar:
	\begin{spverbatim}
	# ./arc2_exp2_script.sh
	# ./arc2_exp_analise.sh 2
	\end{spverbatim}
\noindent
\newline
Para executar os testes da terceira parte, foi gerado o script \verb|arc2_exp3_script.sh|. Seguem os comandos para executar
	\begin{spverbatim}
	# ./arc2_exp3_script.sh
	#./arc2_exp_analise.sh 3
	\end{spverbatim}

\end{document}